% TOP_PROBLEM: {"id": 3380, "title": "The 3n+1 conjecture", "category_id": 1, "category": {"id": 1, "name": "number_theory", "display_name": "Number Theory", "description": "Properties of integers, prime numbers, Diophantine equations.", "slug": "number-theory", "order_index": 1, "created_at": "2026-07-31T15:26:25.670Z"}, "status": "open", "problem_number": "OPG-432", "set_id": 12}
% FIRST_POSTED: 2026-10-01
% ENTRY_KIND: statement_only
% UnsolvedMath ID 3380; source catalogue code OPG-432
% !TeX program = lualatex
% Definition and sources only; this file is not an LLM solution attempt.
\documentclass[11pt,a4paper]{article}
\usepackage[margin=25mm]{geometry}
\usepackage{fontspec}
\setmainfont{lmroman10-regular.otf}[BoldFont=lmroman10-bold.otf,
 ItalicFont=lmroman10-italic.otf,BoldItalicFont=lmroman10-bolditalic.otf]
\usepackage{amsmath,amssymb,mathtools}
\usepackage{unicode-math}
\setmathfont{latinmodern-math.otf}
\AtBeginDocument{\let\setminus\smallsetminus}
\usepackage{xurl,hyperref}
\hypersetup{colorlinks=true,urlcolor=blue}
\setcounter{secnumdepth}{0}
\setlength{\parindent}{0pt}
\setlength{\parskip}{5pt}
\setlength{\emergencystretch}{3em}
\begin{document}
\section*{The 3n+1 conjecture}\label{3380}
{\small UnsolvedMath ID 3380 · OPG-432 · Number Theory · OpenGarden}
\subsection{Definitions and mathematical statement}
Define a map $T:\mathbb Z_{>0}\to\mathbb Z_{>0}$ by
\[
 T(n)=\begin{cases}
 1,&n=1,\\
 n/2,&n>1\text{ even},\\
 3n+1,&n>1\text{ odd}.
 \end{cases}
\]
For an initial positive integer $n_0$, set $n_{j+1}=T(n_j)$ for
$j\ge0$, or equivalently $n_j=T^j(n_0)$. The conjecture asserts
\[
 \forall n_0\in\mathbb Z_{>0}\ \exists j\in\mathbb Z_{\ge0}
                         \quad T^j(n_0)=1.
\]
The convention $T(1)=1$ makes one absorbing. If instead one applies
the odd rule also at one, the equivalent question is eventual entry
into the cycle $1,4,2,1$. No bound on the stopping time $j$ is part
of this statement; it may depend on the starting integer.

Because the iteration is deterministic on positive integers, a
bounded orbit eventually repeats and hence becomes periodic. Thus
a counterexample would either enter a different periodic cycle or
have an unbounded orbit. Establishing absence of additional cycles
alone would leave the unbounded possibility. Conversely, verifying
that every start below a finite cutoff terminates establishes only
that finite range. The problem covers every positive initial integer
with exactly the map above.
\subsection{Short English statement}
Does repeatedly halving an even number and replacing an odd number by three times it plus one always eventually reach one?
\subsection{Sources}
\begin{itemize}
\item[S1] ULAM.AI and the UnsolvedMath Contributors, supplied problems.json, record 3380 (OPG-432), OpenGarden collection. Catalogue identity and source transcription; CC BY 4.0.
\url{https://huggingface.co/datasets/ulamai/UnsolvedMath}.
\item[S2] Open Problem Garden, The 3n+1 conjecture. Original problem and discussion; the mathematical formulation below is expanded from the supplied source record.
\url{http://www.openproblemgarden.org/op/strange_series}.
\end{itemize}
\end{document}
